\chapter[Er trug unsere Krankheit]{Er trug unsere Krankheit}
\quellseite{9}

Die Juden nahmen Christum darum allein nicht an, daß sie ihre eigene Gerechtigkeit ihnen nicht wollten genommen noch gestraft haben, so doch Christus deshalb ins Fleisch kommen und Mensch worden ist, daß wir, die wir Kinder des Zorns und zur Verdammniß verurtheilt waren, durch seine Gerechtigkeit selig und erhalten würden. Darum soll man vornehmlich das Wort: unsere Krankheit und unsere Schmerzen sich wohl einbilden. Denn durch das Wort unser wird uns das selige, heilsame Leiden Christi zugerechnet und mitgetheilt, ja, es wird sogar unser eigen, als hätten wir’s selbst erworben. Denn so er hat getragen meine Schmerzen und meine Krankheit, so ist’s gewiß, daß ich auch ganz frei bin, nicht allein von aller meiner Schuld, sondern auch von meiner wohlverdienten Strafe, und darf deshalb auch nichts fürchten noch verzagen vor dem ernsten Gerichte Gottes. (Luther.)
